\scanpage{177}
% archive sheet number: 176
\mspage{67}
Note that $C_n^v$ is an even function for $n$ even and an odd function for $n$ odd, hence we have

\[
\handwritten{\gamma_{\unsure{m|n}}}=\handwritten{\left\{}
\begin{aligned}
&\frac{w_p}{C_n^{1/2\,p}(1)}\,i\pi\handwritten{\int}\limits_0^1 C^{1/2\,p}(x)(1-x^2)^{1/2\,p-1/2}\,dx && n\text{ odd}\\[4pt]
-&\frac{2w_p}{C_n^{1/2\,p}(1)}\handwritten{\int}\limits_0^1\log x\ C_n^{1/2\,p}(x)(1-x^2)^{1/2\,p-1/2}\,dx && n\text{ even}
\end{aligned}\right.
\]

the result is then proved by evaluating these integrals (see [18a]).

\uline{4.5\quad Reduction of a singular equation to the Fredholm case}.

We now return to the algebra $\handwritten{\mathscr{G}}$ of singular operators on $M$. For each $J\in\handwritten{\mathscr{G}}$ and each $x\in M$ we shall define the symbol of $J$ at $x$; this will be a function on $S_x$. Let $\mathbb{H}_xJ$ be the homogeneous part of $J$ at $x$; we define $\mathbb{S}_xJ=\mathbb{S}\mathbb{H}_xJ$, thus $\mathbb{S}_x$ is the composition of the homomorphism $\mathbb{H}_x$ and the isomorphism $\mathbb{S}$.

\uline{Definition 4.5.1.}\qquad If $J\in\handwritten{\mathscr{G}}$ then the \uline{global symbol} of $J$ is a function on the tangent sphere bundle $S(M)$ defined by

\[
\mathbb{S}_M J(b)=\mathbb{S}_{\pi(b)}J(b).
\]

Note that $\mathbb{S}_M J=0$ implies that $\mathbb{H}_x J=0$ for all $x\in M$ and $J\in\handwritten{\bigcap}\limits_{x\in M}\handwritten{\mathscr{F}}_x=\handwritten{\mathscr{F}}$, so $J$ is completely continuous.
